\section*{結言}
今回の観測では、汎用モデルの性能競争と、推論用アクセラレータやラック設計といった基盤技術の展開が同時に進んだ。法律特化モデルや端末上で動作するマルチモーダルモデルの発表からは、利用場面に合わせた最適化への関心も読み取れる。ただし、ベンチマークの数値や効率改善は主として発表者側の報告である。

また、モデルを実際の業務へ投入する際には、デモの再現性や権限・安全性の管理が引き続き重要となる。Claudeの内部評価事例は二次情報を中心とするため、具体的な経緯の確定には公式資料が必要だ。エージェントのツール利用や検証方法をめぐる議論も、今後の実装と運用を考える材料となる。
